% ===========================================================================
%  PRÉSENTATION DU PROJET — Intégration de MITRE ATT&CK dans un écosystème MCP
%  Partie livrée aujourd'hui : PAGE DE PRÉSENTATION + INTRODUCTION / MISE EN
%  CONTEXTE (diapositives 1 à 6), avec le script complet du présentateur.
%
%  Le script parlé est placé dans les  de chaque diapositive.
%  Pour l'afficher :
%    - sous chaque diapositive :   décommenter  \setbeameroption{show notes}
%    - sur un second écran :       décommenter  la ligne pgfpages ci-dessous
%  Pour le jury : compiler SANS notes (état par défaut).
% ===========================================================================
\documentclass[aspectratio=169,11pt]{beamer}

\usepackage[T1]{fontenc}
\usepackage[utf8]{inputenc}
\usepackage[french]{babel}
\usepackage{lmodern}
\usepackage{tikz}
\usetikzlibrary{arrows.meta,positioning,calc}

% ---------------------------------------------------------------------------
%  Palette (identique au rapport, pour la cohérence visuelle)
% ---------------------------------------------------------------------------
\definecolor{colTitre}{HTML}{1F3A5F}   % bleu profond
\definecolor{colAccent}{HTML}{B5462E}  % brique
\definecolor{colGris}{HTML}{4A4A4A}
\definecolor{colFond}{HTML}{F4F1EA}    % ivoire
\definecolor{colOutil}{HTML}{1F5F4A}   % vert
\definecolor{colOr}{HTML}{C9A227}

% ---------------------------------------------------------------------------
%  Thème sobre et personnalisé
% ---------------------------------------------------------------------------
\setbeamertemplate{navigation symbols}{}
\setbeamertemplate{blocks}[rounded][shadow=false]
\setbeamercolor{structure}{fg=colTitre}
\setbeamercolor{frametitle}{fg=colTitre}
\setbeamercolor{normal text}{fg=colGris}
\setbeamercolor{alerted text}{fg=colAccent}
\setbeamercolor{block title}{bg=colTitre, fg=white}
\setbeamercolor{block body}{bg=colFond, fg=colGris}
\setbeamercolor{block title example}{bg=colOutil, fg=white}
\setbeamercolor{block body example}{bg=colOutil!8, fg=colGris}
\setbeamerfont{frametitle}{series=\bfseries, size=\large}
\setbeamertemplate{itemize item}{\textcolor{colAccent}{\small$\blacksquare$}}
\setbeamertemplate{itemize subitem}{\textcolor{colOr}{\small$\blacktriangleright$}}

% Titre de diapositive avec filet accentué
\setbeamertemplate{frametitle}{%
  \vspace{0.55em}%
  {\usebeamerfont{frametitle}\usebeamercolor[fg]{frametitle}\insertframetitle\par}%
  \vspace{0.18em}{\color{colAccent}\rule{0.30\textwidth}{1.1pt}}\vspace{0.25em}%
}

% Pied de page discret

\newcommand{\grandchiffre}[2]{%
  \begin{center}
    {\fontsize{26}{28}\selectfont\bfseries\color{colAccent}#1}\\[2pt]
    {\small\color{colGris}#2}
  \end{center}}

\title{Intégration du framework MITRE ATT\&CK\\ dans un écosystème MCP pour la détection}
\author{Kamil Saleh}
\date{\today}

\begin{document}



% ===========================================================================
%  DIAPOSITIVE 1 — Page de présentation
% ===========================================================================


\begin{frame}[plain]
  \begin{center}
    \vspace{0.9cm}
    {\color{colAccent}\rule{0.55\textwidth}{1.2pt}}\\[0.55cm]
    {\LARGE\bfseries\color{colTitre}
      Intégration du framework MITRE ATT\&CK\\[0.15cm]
      dans un écosystème MCP pour la détection}\\[0.35cm]
    {\large\itshape\color{colGris}
      Quand l'intelligence artificielle apprend à citer ses sources}\\[0.5cm]
    {\color{colAccent}\rule{0.55\textwidth}{1.2pt}}\\[0.7cm]
    {\normalsize\color{colGris} Kamil Saleh}\\[0.12cm]
    {\small\color{colGris} Université du Québec à Chicoutimi (UQAC)}\\[0.06cm]
    {\small\color{colGris} Projet mené avec la collaboration d'Hydro-Québec & de l'UQAC }\\[0.45cm]
    {\small\color{colGris}\today}
  \end{center}
\end{frame}

% ===========================================================================
%  DIAPOSITIVE 0 — Sommaire
% ===========================================================================

\begin{frame}
  \frametitle{Le chemin de cette présentation}
  \vspace{0.25cm}
  \begin{enumerate}\setlength\itemsep{0.6em}
    \item {\color{colAccent} État de l'art} que sait-on déjà, et que
      manque-t-il ?
    \item Le Model Context Protocol, expliqué simplement
    \item Notre serveur MCP : fonctionnalités, outils, démonstration
    \item Les problèmes rencontrés et corrigés
    \item Ce qui reste à faire : évaluation et sécurisation
  \end{enumerate}
\end{frame}
 

% ===========================================================================
%  DIAPOSITIVE 2 — L'accroche
% ===========================================================================
\begin{frame}
  \begin{center}
    {\Large\bfseries\color{colTitre}
      Suffirait-il de remplacer nos analystes\\ en cybersécurité par l'intelligence artificielle ?}
  \end{center}
  \vspace{0.45cm}
  \begin{itemize}\setlength\itemsep{0.45em}
    \item<2-> Une IA sait \textbf lire et expliquer une attaque mais sait-elle
      \textbf prendre la bonne décision quand la vérité compte ?
    \item<3-> Comment la forcer à s'appuyer sur des faits vérifiables,
      plutôt que sur sa mémoire ?
    \item<4-> C'est exactement le rôle du Model Context Protocol (MCP) et c'est l'objet de ce projet.
  \end{itemize}
\end{frame}

% ===========================================================================
%  DIAPOSITIVE 3 — La réalité d'un SOC
% ===========================================================================
\begin{frame}
  \frametitle{La réalité d'un centre d'opérations de sécurité (SOC)}
  \vspace{0.15cm}
  \begin{columns}[T]
    \begin{column}{0.33\textwidth}
      \grandchiffre{4\,484}{alertes reçues par jour,\\ en moyenne, par équipe SOC\\ {\scriptsize (Vectra, 2023)}}
    \end{column}
    \begin{column}{0.33\textwidth}
      \grandchiffre{67\,\%}{de ces alertes ne sont\\ jamais traitées\\ {\scriptsize (Vectra, 2023)}}
    \end{column}
    \begin{column}{0.33\textwidth}
      \grandchiffre{1 sur 2}{analystes songent à quitter\\ le domaine, épuisés\\ {\scriptsize (Splunk, 2025)}}
    \end{column}
  \end{columns}
  \vspace{0.5cm}
  \begin{block}{ }
    Le danger d'un SOC moderne n'est pas le manque d'information :
    c'est \textbf l'excès d'information non triée et l'épuisement des
    humains qui doivent la trier.
  \end{block}
\end{frame}

% ===========================================================================
%  DIAPOSITIVE 4 — La promesse et le piège
% ===========================================================================
\begin{frame}
  \frametitle{Les LLM : un copilote brillant\ldots{} qui ne dit jamais « je ne sais pas »}
  \vspace{0.2cm}
  \begin{itemize}\setlength\itemsep{0.5em}
    \item Les grands modèles de langage (LLM) savent lire, résumer et
      expliquer une alerte en langage clair une aide précieuse au triage.
    \item<2-> Mais quand l'information leur manque, ils produisent la suite la
      plus \emph{plausible} : c'est l'hallucination .
    \item<3-> En cybersécurité, une technique d'attaque inventée, c'est une enquête orientée vers la mauvaise cible pendant que la vraie
      attaque continue.
  \end{itemize}
  \vspace{0.35cm}
  \begin{exampleblock}{ }<4->
    Le problème n'est pas le LLM lui même mais c'est son ancrage dans les faits.
  \end{exampleblock}
\end{frame}

% ===========================================================================
%  DIAPOSITIVE 5 — Les deux briques : ATT&CK et MCP
% ===========================================================================
\begin{frame}
  \frametitle{Notre réponse : brancher l'IA sur la source officielle}
  \vspace{0.1cm}
  \begin{itemize}\setlength\itemsep{0.35em}
    \item \textbf{MITRE ATT\&CK} : l'encyclopédie mondiale, libre et maintenue,
      des modes opératoires des attaquants « le dictionnaire du cambrioleur ».
    \item \textbf{MCP} \emph{(Model Context Protocol)} : le standard ouvert qui
      branche un modèle d'IA sur une source de données « le port USB-C de l'IA ».
  \end{itemize}
  \vspace{0.3cm}
  \begin{center}
  \begin{tikzpicture}[font=\footnotesize,
    boite/.style={rectangle, rounded corners=3pt, draw, thick, align=center,
      minimum height=1.05cm, inner sep=5pt},
    fleche/.style={-{Stealth[length=2.6mm]}, thick, colGris}]
    \node[boite, draw=colGris, fill=colGris!8] (logs)
      {Alertes\\ et journaux};
    \node[boite, draw=colTitre, fill=colTitre!8, right=1.05cm of logs]
      (agent) {\textbf{Agent (LLM)}\\ raisonne};
    \node[boite, draw=colOutil, fill=colOutil!8, right=1.55cm of agent]
      (serveur) {\textbf{Serveur MCP ATT\&CK}\\ \emph{notre développement}};
    \node[boite, draw=colAccent, fill=colAccent!8, right=1.05cm of serveur]
      (analyste) {Analyste :\\ réponse sourcée};
    \draw[fleche] (logs) -- (agent);
    \draw[fleche, {Stealth[length=2.6mm]}-{Stealth[length=2.6mm]}, colOutil]
      (agent) -- node[above, font=\scriptsize\bfseries, colOutil]{MCP}
      node[below, font=\scriptsize\itshape, colGris]{appels tracés} (serveur);
    \draw[fleche] (serveur -| serveur.east) (analyste);
  \end{tikzpicture}
  \end{center}
  \vspace{0.25cm}
  \begin{block}{ }<2->
    L'agent ne récite plus de mémoire : chaque technique qu'il avance est
    allée chercher, par un appel vérifiable, dans le référentiel officiel.
  \end{block}
\end{frame}

% ===========================================================================
%  DIAPOSITIVE 6 — Le contexte Hydro-Québec et la question de recherche
% ===========================================================================
\begin{frame}
  \frametitle{Pourquoi Hydro-Québec, pourquoi maintenant}
  \vspace{0.15cm}
  \begin{itemize}\setlength\itemsep{0.5em}
    \item Les plateformes commerciales de détection sont des boîtes
      noires : la logique appartient au fournisseur, l'audit est impossible.
    \item Pour une infrastructure critique, la loi canadienne sur la protection des cybersystèmes essentiels (loi C-26) exige une maîtrise technique auditable.
    \item Un serveur MCP développé à l'interne est un actif souverain:
      l'organisation le possède, le lit, le modifie et le journalise.
  \end{itemize}
  \vspace{0.35cm}
\end{frame}


% ===========================================================================
%  DIAPOSITIVE 8 — État de l'art (1/3) : les trois générations du mapping
% ===========================================================================
\begin{frame}
  \frametitle{État de l'art (1/3) - Trois générations pour automatiser le \emph{mapping}}
  \vspace{0.1cm}
  \begin{center}
  \begin{tikzpicture}[
    etape/.style={rectangle, rounded corners=3pt, draw, thick, align=center,
      text width=3.55cm, minimum height=2.15cm, font=\footnotesize, inner sep=4pt},
    fleche/.style={-{Stealth[length=3mm]}, thick, colGris}]
    \node[etape, draw=colTitre, fill=colTitre!8] (g1)
      {\textbf{Classification supervisée}\\ {\scriptsize (ex. TRAM : SciBERT affiné)}\\[3pt]
       {\color{colAccent}\scriptsize $\approx$ 50 techniques couvertes :\\ l'annotation experte plafonne}};
    \node[etape, draw=colOr!80!black, fill=colFond, right=0.95cm of g1] (g2)
      {\textbf{Ancrage par récupération}\\ {\scriptsize (RAG, ex. TechniqueRAG)}\\[3pt]
       {\color{colAccent}\scriptsize jusqu'à $\approx 0{,}90$ de $F_1$, mais\\ les techniques implicites échappent}};
    \node[etape, draw=colOutil, fill=colOutil!8, right=0.95cm of g2] (g3)
      {\textbf{Accès outillé}\\ {\scriptsize (serveurs MCP ATT\&CK)}\\[3pt]
       {\color{colOutil}\scriptsize prometteur mais non encore\\ répliqué indépendamment}};
    \draw[fleche] (g1) -- (g2);
    \draw[fleche] (g2) -- (g3);
    \node[below=1.5mm of g1, font=\scriptsize\itshape, colGris] {$\sim$2021--2024};
    \node[below=1.5mm of g2, font=\scriptsize\itshape, colGris] {2024--2025};
    \node[below=1.5mm of g3, font=\scriptsize\itshape, colGris] {2025--\ldots};
  \end{tikzpicture}
  \end{center}
  \vspace{0.2cm}
  \begin{block}{ }<2->
    Le \emph{mapping} n'est pas de la reconnaissance de mots-clés : c'est un
    raisonnement, qui exige un appui sur une base fiable.
    Chaque génération rapproche le modèle de la source.
  \end{block}
\end{frame}
 
% ===========================================================================
%  DIAPOSITIVE 9 — État de l'art (2/3) : que valent les LLM ?
% ===========================================================================
\begin{frame}
  \frametitle{État de l'art (2/3) - Que valent les LLM face aux attaques ?}
  \vspace{0.2cm}
  \begin{itemize}\setlength\itemsep{0.55em}
    \item Détecter : oui. Ancrés sur des journaux, ils dépassent les
      classifieurs classiques ($F_1 \approx 0{,}93$ contre $\approx 0{,}56$ pour
      XGBoost) et savent expliquer leurs verdicts.
    \item<2-> Se tromper : oui aussi. L'hallucination demeure le risque
      central l'ancrage la \emph{réduit}, il ne l'\emph{élimine} pas.
    \item<3-> Sur le terrain (10 mois d'observation en SOC réel) : les
      analystes utilisent le LLM pour \emph{comprendre}, pas pour \emph{décider}.
  \end{itemize}
  \vspace{0.3cm}
  \begin{exampleblock}{Réponse à notre question d'ouverture}<4->
    La littérature converge : le LLM est une \textbf{augmentation de
    l'analyste}, pas son remplacement.
  \end{exampleblock}
\end{frame}
 
% ===========================================================================
%  DIAPOSITIVE 10 — État de l'art (3/3) : la sécurité du MCP
% ===========================================================================
\begin{frame}
  \frametitle{État de l'art (3/3) - Le revers : la sécurité du MCP lui-même}
  \vspace{0.05cm}
  \begin{itemize}\setlength\itemsep{0.4em}
    \item Un LLM isolé qui se trompe produit du texte ; branché à des
      outils, il peut produire une action.
    \item<2-> Empoisonnement d'outils : 36,5\,\% de réussite moyenne
      sur des serveurs MCP réels et les modèles \emph{les plus capables}
      sont les plus vulnérables (MCPTox, AAAI 2026).
    \item<3-> Vulnérabilités avérées (\texttt{CVE-2025-6514}, gravité 9,6/10) ;
      des garde-fous de la spécification restent optionnels
      (\emph{should} $\neq$ \emph{must}).
    \item<4-> Constat de la NSA (2026) : la prolifération du MCP a
      devancé son modèle de sécurité.
  \end{itemize}
  \vspace{0.1cm}
  \begin{block}{ }<5->
    Un outil d'aide à la détection qui deviendrait lui-même un vecteur
    d'attaque serait une régression d'où le volet sécurisation de
    ce projet.
  \end{block}
\end{frame}
 
% ===========================================================================
%  DIAPOSITIVE 11 — Synthèse de la revue
% ===========================================================================
\begin{frame}
  \frametitle{Synthèse de la revue}
  \begin{columns}[T]
    \begin{column}{0.485\textwidth}
      \begin{exampleblock}{Ce que la littérature établit}
        \footnotesize
        \begin{itemize}\setlength\itemsep{0.35em}
          \item L'ancrage dans une base fiable compte plus que la
            taille du modèle.
          \item Le LLM augmente l'analyste ; il ne le remplace pas.
          \item Les risques du MCP sont réels et documentés.
        \end{itemize}
      \end{exampleblock}
    \end{column}
    \begin{column}{0.485\textwidth}
      \begin{block}{Ce qu'elle laisse ouvert}
        \footnotesize
        \begin{itemize}\setlength\itemsep{0.35em}
          \item Modèles évalués souvent déjà dépassés au moment de la
            publication.
          \item Corpus : du texte CTI en anglais rien sur les
            journaux industriels (OT/ICS).
          \item Benchmarks publics exposés à la contamination ;
            fiabilité résiduelle non caractérisée pour la décision
            critique.
          \item Latence temps réel et sécurité MCP en milieu
            régulé : non évaluées.
        \end{itemize}
      \end{block}
    \end{column}
  \end{columns}
  \vspace{0.12cm}
  \begin{center}
    \onslide<2->{\color{colAccent}\bfseries C'est dans ces angles morts que
    notre projet s'installe  et tout commence par comprendre le MCP.}
  \end{center}
\end{frame}




% ===========================================================================
%  DIAPOSITIVE 12 — Le MCP en une idée
% ===========================================================================
\begin{frame}
  \frametitle{Le MCP en une idée : d'un problème de câbles à une prise universelle}
  \vspace{0.05cm}
  \begin{center}
  \begin{tikzpicture}[
    app/.style={rectangle, rounded corners=2pt, draw=colTitre, thick,
      fill=colTitre!8, minimum width=1.45cm, minimum height=0.52cm, font=\scriptsize},
    src/.style={rectangle, rounded corners=2pt, draw=colOutil, thick,
      fill=colOutil!8, minimum width=1.45cm, minimum height=0.52cm, font=\scriptsize},
    hub/.style={rectangle, rounded corners=3pt, draw=colAccent, very thick,
      fill=colAccent!10, minimum width=1.15cm, minimum height=2.0cm,
      font=\scriptsize\bfseries},
    lien/.style={colGris, thin}]
    % --- Sans standard ---
    \node[app] (a1) at (0,1.2)  {App 1};
    \node[app] (a2) at (0,0.6)  {App 2};
    \node[app] (a3) at (0,0)    {App 3};
    \node[src] (s1) at (3.3,1.2){Source 1};
    \node[src] (s2) at (3.3,0.6){Source 2};
    \node[src] (s3) at (3.3,0)  {Source 3};
    \foreach \i in {1,2,3} \foreach \j in {1,2,3} \draw[lien] (a\i) -- (s\j);
    \node[font=\scriptsize\itshape, colGris, align=center] at (1.65,-0.72)
      {Sans standard : un connecteur\\ sur mesure par couple ($M \times N$)};
    % --- Avec MCP ---
    \begin{scope}[xshift=6.6cm]
    \node[app] (b1) at (0,1.2)  {App 1};
    \node[app] (b2) at (0,0.6)  {App 2};
    \node[app] (b3) at (0,0)    {App 3};
    \node[hub] (m)  at (1.8,0.6) {MCP};
    \node[src] (t1) at (3.6,1.2){Source 1};
    \node[src] (t2) at (3.6,0.6){Source 2};
    \node[src] (t3) at (3.6,0)  {Source 3};
    \foreach \i in {1,2,3} \draw[lien] (b\i) -- (m);
    \foreach \j in {1,2,3} \draw[lien] (m) -- (t\j);
    \node[font=\scriptsize\itshape, colGris, align=center] at (1.8,-0.72)
      {Avec le MCP : une implémentation\\ par côté ($M + N$)};
    \end{scope}
  \end{tikzpicture}
  \end{center}
  \vspace{0.1cm}
  \begin{itemize}\setlength\itemsep{0.3em}
    \item<2-> Standard ouvert : publié par Anthropic (novembre 2024),
      adopté depuis par les grands fournisseurs concurrents.
    \item<3-> Gouvernance neutre : projet confié fin 2025 à une
      fondation sous l'égide de la Linux Foundation personne n'en est
      captif.
  \end{itemize}
\end{frame}
 
% ===========================================================================
%  DIAPOSITIVE 13 — Les trois rôles
% ===========================================================================
\begin{frame}
  \frametitle{Trois rôles et un point de contrôle unique}
  \vspace{0.1cm}
  \begin{center}
  \begin{tikzpicture}[
    comp/.style={rectangle, rounded corners=3pt, draw, thick, align=center,
      font=\footnotesize, inner sep=5pt},
    fleche/.style={{Stealth[length=2.6mm]}-{Stealth[length=2.6mm]}, thick, colGris}]
    \node[comp, draw=colTitre, fill=colTitre!6, minimum width=4.3cm,
      minimum height=2.15cm, label={[font=\footnotesize\bfseries,
      color=colTitre]above:{HÔTE : l'application d'IA}}] (hote) at (0,0) {};
    \node[comp, draw=colTitre, fill=white, minimum width=3.3cm] (llm)
      at (0,0.42) {Modèle (LLM)};
    \node[comp, draw=colTitre, fill=white, minimum width=3.3cm] (client)
      at (0,-0.42) {Client MCP};
    \node[comp, draw=colOutil, fill=colOutil!6, minimum width=4.4cm,
      minimum height=2.15cm, right=2.7cm of hote,
      label={[font=\footnotesize\bfseries, color=colOutil]above:{SERVEUR MCP}}]
      (serveur) {};
    \node[comp, draw=colOutil, fill=white, align=center, minimum width=3.5cm]
      (ref) at (serveur.center) {Référentiel\\ MITRE ATT\&CK};
    \draw[fleche] (client.east) -- node[above, font=\scriptsize, colGris]
      {JSON-RPC 2.0} node[below, font=\scriptsize\itshape, colGris,
      align=center, yshift=-1pt] {connexion dédiée}
      (serveur.west |- client.east);
  \end{tikzpicture}
  \end{center}
  \vspace{0.15cm}
  \begin{block}{ }<2->
    Le serveur ne parle jamais directement au modèle : tout transite
    par le client, donc par l'hôte un guichet unique où chaque appel peut
    être validé, filtré, soumis à confirmation humaine.
  \end{block}
\end{frame}
 
% ===========================================================================
%  DIAPOSITIVE 14 — Les trois primitives
% ===========================================================================
\begin{frame}
  \frametitle{Trois primitives et surtout : qui décide de quoi}
  \vspace{0.2cm}
  \begin{columns}[T]
    \begin{column}{0.32\textwidth}
      \begin{beamercolorbox}[rounded=true, sep=6pt]{block title}
        \centering\textbf{Outils} \emph{(tools)}
      \end{beamercolorbox}
      \begin{beamercolorbox}[rounded=true, sep=6pt]{block body}
        \small Fonctions que le serveur exécute à la demande
        (ex. \texttt{rechercher\_technique}).\\[4pt]
        \textbf{Qui décide : le modèle.}
      \end{beamercolorbox}
    \end{column}
    \begin{column}{0.32\textwidth}
      \begin{beamercolorbox}[rounded=true, sep=6pt]{block title}
        \centering\textbf{Ressources}
      \end{beamercolorbox}
      \begin{beamercolorbox}[rounded=true, sep=6pt]{block body}
        \small Contenus fournis en lecture au contexte
        (ex. une fiche de technique).\\[4pt]
        \textbf{Qui décide : l'application.}
      \end{beamercolorbox}
    \end{column}
    \begin{column}{0.32\textwidth}
      \begin{beamercolorbox}[rounded=true, sep=6pt]{block title}
        \centering\textbf{Prompts}
      \end{beamercolorbox}
      \begin{beamercolorbox}[rounded=true, sep=6pt]{block body}
        \small Modèles d'interaction préconfigurés
        (ex. « analyse cette alerte »).\\[4pt]
        \textbf{Qui décide : l'utilisateur.}
      \end{beamercolorbox}
    \end{column}
  \end{columns}
  \vspace{0.4cm}
  \begin{exampleblock}{ }<2->
    Cette répartition n'est pas un détail technique : c'est un modèle de
    contrôle  elle délimite ce que le modèle peut, et ne peut pas, provoquer.
  \end{exampleblock}
\end{frame}
 
% ===========================================================================
%  DIAPOSITIVE 15 — Le cycle de vie d'un échange
% ===========================================================================
\begin{frame}
  \frametitle{Un échange MCP : trois phases, aucune improvisation}
  \vspace{0.1cm}
  \begin{center}
  \begin{tikzpicture}[
    phase/.style={rectangle, rounded corners=3pt, draw, thick, align=center,
      text width=3.55cm, minimum height=2.05cm, font=\footnotesize, inner sep=4pt},
    fleche/.style={-{Stealth[length=3mm]}, thick, colGris}]
    \node[phase, draw=colTitre, fill=colTitre!8] (p1)
      {\textbf{1 $\cdot$ Initialisation}\\[3pt]
       \scriptsize On s'accorde : version du\\ protocole + capacités de chacun.\\
       {\color{colAccent}Rien de non négocié ne\\ pourra être invoqué.}};
    \node[phase, draw=colOutil, fill=colOutil!8, right=0.95cm of p1] (p2)
      {\textbf{2 $\cdot$ Opération}\\[3pt]
       \scriptsize \texttt{tools/list} : le catalogue.\\
       \texttt{tools/call} : l'invocation.\\
       {\color{colOutil}Chaque résultat = un\\ message tracé.}};
    \node[phase, draw=colAccent, fill=colAccent!8, right=0.95cm of p2] (p3)
      {\textbf{3 $\cdot$ Arrêt}\\[3pt]
       \scriptsize Fermeture du canal\\ de transport.\\
       \scriptsize\itshape (pas de message dédié)};
    \draw[fleche] (p1) -- (p2);
    \draw[fleche] (p2) -- (p3);
  \end{tikzpicture}
  \end{center}
  \vspace{0.25cm}
  \begin{block}{L'idée essentielle}<2->
    Le modèle ne fabrique pas les données qu'il présente : il les
    obtient par des messages identifiés et journalisables. C'est la
    propriété qui motive tout notre projet.
  \end{block}
\end{frame}
 
% ===========================================================================
%  DIAPOSITIVE 16 — Les deux transports
% ===========================================================================
\begin{frame}
  \frametitle{Deux façons de se parler : \texttt{stdio} et Streamable HTTP}
  \vspace{0.1cm}
  \begin{columns}[T]
    \begin{column}{0.485\textwidth}
      \begin{exampleblock}{\texttt{stdio} la même pièce}
        \footnotesize
        \begin{itemize}\setlength\itemsep{0.3em}
          \item Le serveur est un sous-processus local du client.
          \item Aucune exposition réseau ; cycle de vie adossé au
            processus parent.
          \item Cas d'usage : poste d'analyste, développement.
        \end{itemize}
      \end{exampleblock}
    \end{column}
    \begin{column}{0.485\textwidth}
      \begin{block}{Streamable HTTP l'appel réseau}
        \footnotesize
        \begin{itemize}\setlength\itemsep{0.3em}
          \item Le serveur est un service réseau : un point d'entrée,
            plusieurs clients.
          \item Exposition réelle : certaines protections exigées,
            d'autres seulement \emph{recommandées}.
          \item Cas d'usage : intégration à l'infrastructure d'un SOC.
        \end{itemize}
      \end{block}
    \end{column}
  \end{columns}
  \vspace{0.25cm}
  \begin{itemize}
    \item<2-> Même format de message dans les deux cas : changer de
      transport ne change ni les outils, ni leur logique mais change la
      surface d'attaque (chapitre sécurisation).
    \item<3-> Protocole en évolution rapide : révision stable
      2025-11-25 ; une révision 2026-07-28 en cours de finalisation.
  \end{itemize}
\end{frame}
 
% ===========================================================================
%  DIAPOSITIVE 17 — Périmètre : ce que le MCP apporte, ce qu'il ne fait pas
% ===========================================================================
\begin{frame}
  \frametitle{Soyons précis : ce que le MCP apporte et ce qu'il ne fait pas}
  \vspace{0.1cm}
  \begin{columns}[T]
    \begin{column}{0.485\textwidth}
      \begin{exampleblock}{Ce qu'il apporte au projet}
        \footnotesize
        \begin{itemize}\setlength\itemsep{0.3em}
          \item \textbf{Interopérabilité} : un serveur conforme se branche sur
            tout hôte conforme.
          \item \textbf{Ancrage vérifiable} : un identifiant retourné provient
            nécessairement du référentiel.
          \item \textbf{Traçabilité} : chaque appel est journalisable 
            l'auditabilité qu'exige la LPCE (ex-C-26).
          \item \textbf{Souveraineté} : la logique de détection appartient à
            l'organisation.
        \end{itemize}
      \end{exampleblock}
    \end{column}
    \begin{column}{0.485\textwidth}
      \begin{block}{Ce qu'il ne fait pas}
        \footnotesize
        \begin{itemize}\setlength\itemsep{0.3em}
          \item Il ne rend pas le modèle plus intelligent : il
            normalise l'\emph{échange de contexte}, pas le raisonnement.
          \item Il ne garantit pas la qualité du mapping
            c'est la conception de notre serveur.
          \item Il ne garantit pas la sécurité applicative c'est
            notre volet sécurisation.
        \end{itemize}
      \end{block}
    \end{column}
  \end{columns}
  \vspace{0.2cm}
  \begin{center}
    \onslide<2->{\color{colAccent}\bfseries Le MCP fournit le canal pas la
    méthode. Voyons maintenant le serveur que nous avons construit dessus.}
  \end{center}
\end{frame}
 
% ===========================================================================
%  DIAPOSITIVE 18 — Notre serveur : la carte d'identité
% ===========================================================================
\begin{frame}
  \frametitle{Notre serveur MCP : la carte d'identité}
  \vspace{0.05cm}
  \begin{columns}[T]
    \begin{column}{0.485\textwidth}
      \begin{exampleblock}{Ce qu'il est}
        \footnotesize
        \begin{itemize}\setlength\itemsep{0.3em}
          \item \textbf{Un fichier Python unique}, bibliothèque standard
            seulement auditable ligne par ligne, zéro dépendance tierce.
          \item Transport stdio local (mode développement) 
            l'exposition réseau attend le volet sécurisation.
          \item \textbf{Lecture seule} : aucun outil d'écriture le modèle
            ne peut pas polluer la source de vérité.
        \end{itemize}
      \end{exampleblock}
    \end{column}
    \begin{column}{0.485\textwidth}
      \begin{block}{Ce qu'il sert}
        \footnotesize
        \begin{itemize}\setlength\itemsep{0.3em}
          \item Base hybride Enterprise + ICS (ATT\&CK 19.1) :
            27 tactiques, 794 techniques, 176 groupes, 827 logiciels,
            96 mitigations.
          \item 14 outils en 4 familles : cycle de mapping,
            consultation, référentiel interne, exploitation.
          \item Chaque réponse datée (\texttt{db\_version}) + journal
            d'audit JSONL.
        \end{itemize}
      \end{block}
    \end{column}
  \end{columns}
  \vspace{0.25cm}
  \begin{center}
    \onslide<2->{\color{colTitre}\bfseries Chaque ligne de cette carte
    d'identité traduit une exigence de la problématique en décision technique
    vérifiable.}
  \end{center}
\end{frame}
 
% ===========================================================================
%  DIAPOSITIVE 19 — Architecture interne
% ===========================================================================
\begin{frame}
  \frametitle{L'architecture : le trajet d'une requête}
  \vspace{0.05cm}
  \begin{center}
  \begin{tikzpicture}[font=\scriptsize,
    blk/.style={rectangle, rounded corners=2.5pt, draw, thick, align=center,
      minimum height=0.85cm, text width=1.85cm, inner sep=3pt},
    io/.style={rectangle, rounded corners=2.5pt, draw=colGris!70, fill=colGris!8,
      align=center, minimum height=0.7cm, text width=1.3cm, inner sep=3pt},
    dbb/.style={rectangle, rounded corners=2.5pt, draw=colOr!80!black,
      fill=colFond, align=center, text width=2.1cm, inner sep=3pt,
      minimum height=0.8cm},
    fl/.style={-{Stealth[length=2.4mm]}, thick, colGris}]
    \node[io] (in) {\texttt{stdin}\\ JSON-RPC};
    \node[blk, draw=colTitre, fill=colTitre!8, right=0.5cm of in] (rout)
      {\textbf{Routage}\\ décodage};
    \node[blk, draw=colTitre, fill=colTitre!8, right=0.5cm of rout] (val)
      {\textbf{Validation}\\ types, longueurs};
    \node[blk, draw=colOutil, fill=colOutil!8, right=0.5cm of val] (mot)
      {\textbf{Moteur}\\ recherche, scoring};
    \node[blk, draw=colAccent, fill=colAccent!8, right=0.5cm of mot] (ser)
      {\textbf{Sérialisation}\\ + \texttt{db\_version}};
    \node[io, right=0.5cm of ser] (out) {\texttt{stdout}\\ JSON-RPC};
    \draw[fl] (in)--(rout); \draw[fl] (rout)--(val); \draw[fl] (val)--(mot);
    \draw[fl] (mot)--(ser); \draw[fl] (ser)--(out);
    \node[dbb, below=0.55cm of mot] (base)
      {Base ATT\&CK E+ICS\\ + index lexicaux};
    \node[dbb, left=0.45cm of base, text width=1.8cm] (thq)
      {Référentiel\\ interne (THQ)};
    \node[dbb, below=0.55cm of ser, text width=1.8cm, draw=colGris!70] (aud)
      {Journal d'audit\\ JSONL};
    \draw[{Stealth[length=2.4mm]}-{Stealth[length=2.4mm]}, colGris!70] (mot)--(base);
    \draw[{Stealth[length=2.4mm]}-{Stealth[length=2.4mm]}, colGris!70]
      (mot.south west)--(thq.north);
    \draw[fl, colGris!70] (ser)--(aud);
  \end{tikzpicture}
  \end{center}
  \vspace{0.15cm}
  \begin{block}{L'idée essentielle}<2->
    La validation précède toujours le traitement, chaque réponse est
    datée, chaque appel laisse une trace et le moteur est
    \textbf{indépendant du transport} : passer en réseau n'ajoutera qu'une
    couche en amont, sans toucher au moteur ni aux réponses.
  \end{block}
\end{frame}
 
% ===========================================================================
%  DIAPOSITIVE 20 — La méthode : proposer puis vérifier
% ===========================================================================
\begin{frame}
  \frametitle{La méthode encodée : proposer, puis vérifier}
  \vspace{0.15cm}
  \begin{itemize}\setlength\itemsep{0.45em}
    \item Proposer \texttt{search\_techniques} retourne des
      \emph{candidats} : score normalisé $[0..1]$ ($\geq 0{,}6$ fort ;
      $< 0{,}3$ faible) et preuves les mots réellement trouvés.
    \item<2-> Vérifier \texttt{get\_technique} confirme chaque
      candidat retenu : définition complète, tactiques, détections,
      sous-techniques.
    \item<3-> Règles d'honnêteté  la technique parente d'abord ;
      et « pas de mapping » est une \emph{conclusion valide} : le serveur
      l'écrit lui-même.
  \end{itemize}
  \vspace{0.3cm}
  \begin{block}{L'idée essentielle}<4->
    Cette discipline n'est pas demandée poliment au modèle : elle est
    \textbf{écrite dans la description des outils} elle voyage avec le
    serveur, versionnée avec le code, identique pour tous les modèles.
  \end{block}
\end{frame}
 
% ===========================================================================
%  DIAPOSITIVE 21 — Le référentiel interne THQ
% ===========================================================================
\begin{frame}
  \frametitle{Le référentiel interne THQ : étendre ATT\&CK sans le trahir}
  \vspace{0.15cm}
  \begin{itemize}\setlength\itemsep{0.45em}
    \item Des fiches \texttt{THQ}$xxxx$ pour les comportements propres à
      l'organisation toujours rattachées à une technique MITRE.
    \item<2-> MITRE d'abord, l'interne ensuite : recherche anti-doublon
      soumise aux mêmes règles de score et de preuves.
    \item<3-> Lecture seule par conception : le modèle \emph{propose}
      un brouillon ; un humain relit, valide et committe le fichier.
  \end{itemize}
  \vspace{0.3cm}
  \begin{exampleblock}{La boucle de renseignement locale}<4->
    Incident interne $\rightarrow$ fiche brouillon proposée $\rightarrow$
    validation humaine $\rightarrow$ le prochain analyste la retrouve,
    accrochée à la technique MITRE qu'elle précise
    (\texttt{internal\_extensions}).
  \end{exampleblock}
\end{frame}
 
% ===========================================================================
%  DIAPOSITIVE 22 — La trajectoire de versionnement
% ===========================================================================
\begin{frame}
  \frametitle{Six versions, chacune née d'un défaut constaté à l'usage}
  \vspace{0.2cm}
  \begin{center}
  \begin{tikzpicture}[font=\scriptsize]
    \draw[colGris!50, line width=1pt] (0,0) -- (12.6,0);
    \foreach \x/\v in {0.5/v1.0, 2.9/v1.1, 5.3/v2.0, 7.7/v2.1, 10.1/v3.0, 12.1/v3.1}{
      \filldraw[colTitre] (\x,0) circle (2.2pt);
      \node[colTitre, font=\scriptsize\bfseries, above=2pt] at (\x,0) {\v};}
    \node[text width=2.1cm, align=center, below=4pt, colGris] at (0.5,0)
      {Enterprise seul\\ \itshape point de départ};
    \node[text width=2.1cm, align=center, below=4pt, colGris] at (2.9,0)
      {\textbf{convergence IT/OT}\\ journal d'audit};
    \node[text width=2.2cm, align=center, below=4pt, colGris] at (5.3,0)
      {\textbf{refonte du moteur}\\ score, détection};
    \node[text width=2.0cm, align=center, below=4pt, colGris] at (7.7,0)
      {\textbf{référentiel interne}\\ 14 outils};
    \node[text width=2.0cm, align=center, below=4pt, colGris] at (10.1,0)
      {\textbf{jalon production}\\ sans état};
    \node[text width=2.0cm, align=center, below=4pt, colGris] at (12.1,0)
      {\textbf{jalon auditabilité}};
  \end{tikzpicture}
  \end{center}
  \vspace{0.45cm}
  \begin{alertblock}{Le fil rouge de tout ce qui suit}<2->
    Les défauts les plus graves n'ont jamais produit d'erreur : ils se
    présentaient comme des réponses normales. La \emph{défaillance silencieuse}
    est le mode d'échec caractéristique d'un composant d'ancrage.
  \end{alertblock}
\end{frame}
 
% ===========================================================================
%  DIAPOSITIVE 23 — Défaut 1 : la couverture IT/OT
% ===========================================================================
\begin{frame}
  \frametitle{Défaut 1 : Un serveur qui ignorait la moitié du monde}
  \vspace{0.1cm}
  \begin{columns}[T]
    \begin{column}{0.485\textwidth}
      \begin{alertblock}{Avant (v1.0)}
        \footnotesize Matrice Enterprise seule. Sur des journaux industriels
        (Modbus, automates) : aucun candidat pertinent ou pire,
        des techniques bureautiques superficiellement voisines.
      \end{alertblock}
    \end{column}
    \begin{column}{0.485\textwidth}
      \begin{exampleblock}{Après (v1.1)}<2->
        \footnotesize Fusion \textbf{Enterprise + ICS} : chaque objet étiqueté
        par sa matrice, rattachements cloisonnés « Persistence »
        existe deux fois (\texttt{TA0003} $\neq$ \texttt{TA0110}), jamais
        confondues.
      \end{exampleblock}
    \end{column}
  \end{columns}
  \vspace{0.35cm}
  \begin{block}{L'idée essentielle}<3->
    Le défaut n'était pas algorithmique, il était \textbf{ontologique} : la
    base interrogée ne couvrait pas le monde observé exactement l'angle
    mort OT identifié dans l'état de l'art.
  \end{block}
\end{frame}
 
% ===========================================================================
%  DIAPOSITIVE 24 — Défaut 2 : recherche et score
% ===========================================================================
\begin{frame}
  \frametitle{Défaut 2 : Une recherche qui trouvait trop, un score qui ne disait rien}
  \vspace{0.1cm}
  \begin{columns}[T]
    \begin{column}{0.485\textwidth}
      \begin{alertblock}{Avant (v1.1)}
        \footnotesize \texttt{"port"} $\rightarrow$ \textbf{264 candidats},
        dont \emph{Security Sup\textbf{port} Provider} et
        \emph{\textbf{Port}able Executable Injection}.\\[3pt]
        Score : un entier sans échelle « 8 »\ldots{} sur quoi ?
      \end{alertblock}
    \end{column}
    \begin{column}{0.485\textwidth}
      \begin{exampleblock}{Après (v2.0)}<2->
        \footnotesize Mots \textbf{entiers} + racinisation :
        \texttt{"port"} $\rightarrow$ \textbf{37 candidats, zéro faux
        positif}.\\[3pt]
        Score normalisé $[0..1]$, seuils documentés (0,6 / 0,3), et
        \textbf{preuves} : les mots effectivement trouvés.
      \end{exampleblock}
    \end{column}
  \end{columns}
  \vspace{0.35cm}
  \begin{block}{L'idée essentielle}<3->
    Un candidat n'est utile que s'il est \textbf{interprétable} : un score qui
    a une échelle, et la preuve de ce qui a réellement correspondu.
  \end{block}
\end{frame}
 
% ===========================================================================
%  DIAPOSITIVE 25 — Défaut 3 : le référentiel est vivant
% ===========================================================================
\begin{frame}
  \frametitle{Défaut 3 : Le référentiel est vivant : deux pièges silencieux}
  \vspace{0.1cm}
  \begin{itemize}\setlength\itemsep{0.5em}
    \item \textbf{MITRE a restructuré son modèle de détection} (v18) : nos
      champs se sont vidés sur \emph{toutes} les techniques sans une seule
      erreur. \onslide<2->{Corrigé : lecture du nouveau modèle
      (stratégies $\rightarrow$ analytics $\rightarrow$ sources de logs)
      \textbf{0 $\rightarrow$ 387 sources de données}, 794/794 techniques
      dotées de détections.}
    \item<3-> \textbf{MITRE révoque des techniques} (158 redirections en 19.1) :
      \texttt{T1086} répondait « non trouvée » à un rapport CTI de 2019.
      \onslide<4->{Corrigé : redirection explicite
      \texttt{T1086} $\rightarrow$ « révoquée, consulter \texttt{T1059.001} ».}
  \end{itemize}
  \vspace{0.25cm}
  \begin{block}{L'idée essentielle}<5->
    S'ancrer sur un référentiel vivant impose de suivre sa vie :
    versions datées dans chaque réponse, redirections plutôt que silences.
  \end{block}
\end{frame}
 
% ===========================================================================
%  DIAPOSITIVE 26 — Défaut 4 : le journal qui ne prouvait pas
% ===========================================================================
\begin{frame}
  \frametitle{Défaut 4 : Un journal qui ne prouvait pas ce qu'il devait prouver}
  \vspace{0.1cm}
  \begin{itemize}\setlength\itemsep{0.45em}
    \item Une session de \textbf{6 messages} de protocole ne laissait
      qu'\textbf{une} ligne d'audit ; les erreurs, elles, n'en laissaient
      aucune.
    \item<2-> Corrigé (v3.1) : \textbf{chaque phase journalisée} (1
      $\rightarrow$ 6 événements) y compris l'appel d'un outil
      \emph{inexistant} : l'\textbf{hallucination d'outil devient mesurable}
      depuis le seul journal.
    \item<3-> Enregistrements \textbf{enrichis} : identifiant de requête,
      latence, version du référentiel, résumé des arguments et du résultat
      et les pertes de traces \textbf{se comptent elles-mêmes}.
  \end{itemize}
  \vspace{0.25cm}
  \begin{block}{L'idée essentielle}<4->
    Un journal destiné à servir de preuve doit pouvoir déclarer ses
    propres lacunes : sinon l'absence d'un événement devient indiscernable
    de sa non-occurrence.
  \end{block}
\end{frame}
 
% ===========================================================================
%  DIAPOSITIVE 27 — Bilan du serveur et transition
% ===========================================================================
\begin{frame}
  \frametitle{Bilan : ce que le serveur garantit désormais}
  \vspace{0.1cm}
  \begin{columns}[T]
    \begin{column}{0.485\textwidth}
      \begin{exampleblock}{Garanties acquises}
        \footnotesize
        \begin{itemize}\setlength\itemsep{0.3em}
          \item Couverture \textbf{IT/OT} cloisonnée par matrice.
          \item Candidats \textbf{interprétables} : score étalonné + preuves.
          \item Référentiel vivant \textbf{suivi} : versions datées,
            redirections.
          \item Traçabilité \textbf{de bout en bout}, pertes comptées.
        \end{itemize}
      \end{exampleblock}
    \end{column}
    \begin{column}{0.485\textwidth}
      \begin{block}{La leçon de méthode}
        \footnotesize La défaillance silencieuse étant le mode d'échec
        caractéristique d'un composant d'ancrage, chaque correction est
        démontrée par un protocole rejouable, la même requête aux deux
        versions, les deux sorties comparées jamais seulement affirmée.
      \end{block}
    \end{column}
  \end{columns}
  \vspace{0.35cm}
  \begin{center}
    \onslide<2->{\color{colAccent}\bfseries Le serveur existe et il est
    honnête. Reste à le \emph{prouver} l'évaluation et à le
    \emph{durcir}  la sécurisation.}
  \end{center}
\end{frame}
 


% ===========================================================================
%  SUITE DU PLAN (à construire lors des prochaines séances) :
%    - Présentation du sommaire
%    - État de l'art
%    - Introduction au MCP
%    - Notre serveur MCP (fonctionnalités, outils, fonctionnement, démonstration)
%    - Les problèmes corrigés (explications techniques + démonstrations)
%    - Ce qui reste à faire (évaluation + sécurisation)
%    - L'évaluation (méthodologie)
%    - La sécurisation du MCP (transports, journalisation, authentification,
%      protection contre les attaques)
% ===========================================================================

\end{document}
